\documentclass{article}
\usepackage[utf8]{inputenc}
\usepackage{amsmath}
\usepackage{geometry}
\geometry{margin=2cm}
\begin{document}

\section*{Questão ENEM 2024 - Ciências da Natureza: Geração de Eletricidade com Grafeno e Água da Chuva}

\subsection*{Interpretação do Enunciado e Estratégia de Resposta}

A questão apresenta uma inovação tecnológica que utiliza uma fina camada de grafeno sobre placas solares para possibilitar a geração de eletricidade mesmo em dias chuvosos. Isso ocorre pela interação entre os elétrons livres do grafeno e os íons dissociados da água da chuva (Am+ e Bn-). O objetivo é compreender o fenômeno físico que possibilita essa geração energética.

Nossa estratégia é explorar o enunciado e a imagem que evidencia camadas finas de água contendo íons e camada de grafeno contendo elétrons, interpretando como essa interação entre prótons (Am+) e elétrons (e-) causa um movimento desses elétrons e, consequentemente, uma diferença de potencial elétrico.

\bigskip

\subsection*{Revisão Conceitual e Conteúdos Essenciais}

\begin{tabular}{p{6cm} p{10cm}}
\textbf{Conceito} & \textbf{Descrição e Aplicação} \\
\hline
Íons e Movimento de Elétrons & Íons positivos Am+ atraem elétrons (cargas negativas), induzindo o movimento destes na superfície condutora do grafeno. Este movimento de cargas gera corrente elétrica.\\
Diferença de Potencial Elétrico & O movimento desigual de cargas elétricas gera uma voltagem (diferença de potencial) entre regiões diferentes, fundamento básico de geração de energia elétrica em dispositivos.\\
\end{tabular}

\bigskip

\noindent O grafeno consiste numa monocamada de carbono com alta condutividade eletrônica, permitindo o livre movimento dos elétrons. A interação eletrostática entre os íons positivos da água da chuva e os elétrons livres do grafeno provoca a formação da diferença de potencial.

\subsection*{Resolução e "Pulo do Gato"}

Ao analisar a imagem e o enunciado, identificamos que a camada fina de água contendo os íons Am+ atrai os elétrons da camada de grafeno abaixo dela, deslocando-os. Esse deslocamento gera regiões de alto e baixo potencial elétrico, criando uma voltagem (diferença de potencial) que pode ser aproveitada para gerar eletricidade.

\medskip

\noindent \textbf{Pulo do Gato:} A geração de eletricidade aqui é um fenômeno físico fundamentalmente baseado no movimento e separação das cargas elétricas (elétrons) induzida pela presença dos íons positivos, e não uma reação química.

\subsection*{Armadilhas e Distratores}

\begin{tabular}{|c|p{5cm}|p{3cm}|p{5cm}|p{5cm}|}
\hline
\textbf{Opção} & \textbf{Raciocínio Possível} & \textbf{Tipo de Equívoco} & \textbf{Por Que Parece Plausível} & \textbf{Por Que Está Errado} \\
\hline
A & Supor que há oxidação química dos íons pelo grafeno & Confusão sobre reação química & Íons e grafeno sugerem reações químicas & Fenômeno é puramente elétrico; não há oxidação no processo \\
\hline
B & Entender o grafeno como barreira física da água & Má interpretação do papel do grafeno & Grafeno confere barreira física em outros contextos & Aqui o foco é geração de eletricidade, não barreira à água \\
\hline
C & Achar que o grafeno serve como catalisador diminuindo energia de ativação & Aplicação indevida de conceitos catalíticos & Grafeno conhecido por atuar cataliticamente em outras reações & O processo no enunciado é fenômeno físico, não reação química catalisada \\
\hline
D & Imaginar formação de compósito químico entre grafeno e íons & Erro conceitual em química de materiais & Combinações sugerem novos compósitos & Ocorre apenas interação eletrostática, sem formação de nova substância \\
\hline
\end{tabular}

\subsection*{Código da Questão}
\texttt{CN\_ELETRICIDADE\_PF\_001}

\bigskip

\noindent Esta resolução apresenta uma análise coerente entre enunciado, imagem e alternativas, indicando claramente a alternativa E como correta devido à explicação física do fenômeno de geração de eletricidade com grafeno e água da chuva.

\end{document}